% !TeX root = ../main.tex

\chapter{Analyse, voorspelling en strategie}
\label{chap:analyse}

De lokaal opgeslagen timingdata wordt door de analyselaag omgezet in informatie die de race-engineer tijdens een sessie kan gebruiken. Daarbij worden alleen representatieve rondes en sectoren meegenomen. Op basis daarvan berekent de applicatie statistieken, voorspellingen, klassevergelijkingen en pitstopprojecties die strategische beslissingen ondersteunen.

\section{Ronde- en pilootstatistieken}
Voor een reeks geldige ronden worden laatste tijd, beste tijd, gemiddelde tijd, sectorgemiddelden en
beste sectoren bepaald. Pilootstatistieken worden per combinatie van wagen en piloot opgebouwd. De
beste piloot van het eigen team is niet hardcoded: wanneer de actuele piloot na zijn stint
sneller blijkt, kan hij de vorige referentiepiloot vervangen.

In racemodus vergelijkt het dashboard onder meer:

\begin{itemize}
  \item beste teamtijd met de laatste ronde van de huidige piloot
  \item beste teamtijd met de beste ronde van de huidige piloot
  \item gemiddeld tempo van de beste piloot met dat van de huidige piloot
  \item tempo van de beste wagen in de klasse met de eigen actuele stint
  \item tempo van een vrij gekozen klassewagen met de eigen actuele stint
\end{itemize}

In qualifyingmodus zijn gemiddelden minder relevant. Daar gebruikt de app vooral beste en laatste
geldige ronde van teamgenoten en concurrenten. Practice legt de nadruk op referenties en consistente
vergelijking zonder pitstopstrategie.

\section{Voorspelling van de actuele ronde}
De rondevoorspelling is bewust eenvoudig en uitlegbaar. Voor de huidige piloot worden per sector
maximaal de laatste 5 geldige sectortijden geselecteerd. Na sector 1 is de voorspelling:

\[ \hat{T}_{lap} = T_{S1,actueel} + \overline{T}_{S2,recent} + \overline{T}_{S3,recent}. \]

Na sector 2 worden beide actuele sectoren gebruikt:

\[ \hat{T}_{lap} = T_{S1,actueel} + T_{S2,actueel} + \overline{T}_{S3,recent}. \]

Na sector 3 wordt de voorspelling niet vervangen door de echte rondetijd. Zo blijft zichtbaar wat na
sector 2 voorspeld was. Tussen de finishlijn en de nieuwe sector 1 toont het dashboard de afwijking
tussen voorspelling en gerealiseerde ronde. Dit geeft onmiddellijke feedback over de kwaliteit van
de voorspellingen.

Voor een nieuwe piloot is na zijn eerste geldige volledige ronde voldoende basisdata beschikbaar om
vanaf sector 1 van zijn tweede ronde een eerste voorspelling te maken. Wanneer een piloot eerder al
een stint reed, wordt zijn eigen historische sectordata opnieuw gebruikt.

\section{Referentietijden}
De gebruiker kan een referentierondetijd en drie sectorreferenties ingeven via editknoppen in de vakken. De waarden worden per sessiemodus opgeslagen. \code{normReference.js} berekent de
delta en deelt de toestand in als veilig, nabij of kritisch. Deze grenswaarden kunnen ook aangepast worden.
Het vak voor predicted lap time krijgt een duidelijke groene, oranje of rode achtergrond. 
De ideal
time is de som van de beste sectoren van de wagen en krijgt
eveneens een status. Bij elke sectorreferentie worden twee marges getoond: die van de laatste sector
en die van de beste sector. Hierdoor ziet de engineer zowel de actuele uitvoering als het
theoretische risico.

\section{Vergelijkingen binnen de klasse}
Een tijdsverschil tussen twee klassewagens mag niet worden afgeleid door alleen hun zichtbare
klassepositie te vergelijken. Er kunnen wagens uit andere klassen tussen staan. 
Wanneer wagens in dezelfde ronde zitten, combineert de applicatie de actuele gap met het verschil
tussen recente gemiddelden. Zo ontstaat een indicatie van het aantal ronden of de tijd tot een
mogelijke inhaalactie. 

\section{Pitstopplanner}
De pitstopplanner gebruikt configureerbare regels voor totale raceduur, gesloten pitwindow tijd aan de start en einde van de race, cooldown na een geldige stop, aantal verplichte stops en geschatte pitduur. De balk
toont rode en groene zones en bewaart reeds verstreken rode cooldownsegmenten. Alleen stops die
binnen een geldig groen venster plaatsvinden, tellen mee voor het verplichte aantal.

De planner berekent ook het laatste veilige moment waarop nog voldoende verplichte stops met
cooldown kunnen plaatsvinden. De pitduur blijft tijdens de race aanpasbaar,
omdat een bandenwissel en een pilotenwissel verschillende tijden kunnen vragen.

Voor de projectie na een pitstop worden de relatieve verschillen in de algemene rangschikking
opgebouwd totdat de geschatte pit loss is overschreden. Vervolgens wordt bepaald tussen welke
klassewagens de eigen wagen zou uitkomen.

\section{Full Course Yellow}
Onder FCY rijdt het veld met een opgelegde snelheid en kan de relatieve pit loss kleiner zijn.
Circuitprofielen bevatten daarom de afstand die een niet-pittende wagen tussen pit entry en pit exit
over het gewone circuit aflegt en de FCY-snelheid. De tijd op dat traject is:

\[ T_{track,FCY} = \frac{d_{track}}{v_{FCY}/3{,}6}. \]

De geschatte netto pit loss vergelijkt deze trajecttijd en de gap met de ingestelde pitduur. Hierdoor kan er ook tijdens een FCY een juiste inschatting gemaakt worden van de juiste pit loss en de positie in het veld na een pitstop. 

\section{Automatische rapportage}
Het idee om automatische PDF-rapporten toe te voegen ontstond tijdens de praktijktest in Spa. Na
zijn stint vroeg een piloot hoe zijn stint verlopen was. Op dat moment kon ik vooral zeggen dat hij
ronde na ronde iets sneller werd, maar ik had nog geen duidelijk overzicht van waar die verbetering
vandaan kwam: welke sectoren beter werden, welke ronden representatief waren, hoe constant de stint
was en hoe hij het deed in vergelijking met teamgenoten of klassewagens. Ook Jan vroeg na de race om een
verslag van de volledige wedstrijd. Dat maakte duidelijk dat deze rapporten niet alleen nuttig zijn
voor onmiddellijke feedback aan de piloot, maar ook voor latere evaluatie.

Naast het live dashboard kan de applicatie daarom PDF-rapporten genereren op basis van de lokaal
opgeslagen racehistoriek. Per pilootstint wordt een apart document gemaakt met de rondetijden,
sectoren, sectorgemiddelden, beste en traagste geldige ronden, uitgesloten ronden en vergelijkingen met teamgenoten en klassewagens. De rapporten gebruiken
dezelfde centrale selectieregels als het dashboard: FCY-, Safety-Car-, pit-in-, outlap- en
duidelijke outlierronden blijven bewaard, maar worden niet gebruikt voor gemiddelden te berekenen.

Voor de volledige race wordt ook een overzichtspagina gegenereerd. Deze bundelt onder
andere het aantal gereden ronden, geldige ronden, gemiddelde rondetijd, beste ronde,
sectorgemiddelden, pitstopinformatie en vergelijking tussen piloten. Hierdoor kan
de engineer na de sessie sneller zien waar tijd werd gewonnen of verloren, welke stint het meest
constant was en hoe sterk de pitstops afweken van de verwachte pitduur. De PDF's zijn dus geen
aparte databron, maar een leesbare samenvatting van dezelfde JSON-, JSONL- en CSV-bestanden die
tijdens de race werden opgebouwd. Een voorbeeld van een automatisch gegenereerde race-PDF op basis
van de Spa-data is opgenomen in de appendix in Bijlage~\ref{app:spa-race-pdf}.
